\section{Conclusions}

\begin{frame}{Conclusions}
	Dans ce cours, nous avons vu:
	\begin{itemize}
		\item Que l'analyse en composantes principales est une solution \textbf{optimale} au problème de combinaison \textbf{linéaire} des caractéristiques d'entrée afin de maximiser la variance.
		\pause
		\item Que l'analyse en composantes principales permet de \textbf{redistribuer} la variance des caractéristiques \textbf{sans perte}, tout en rendant les combinaisons linéaires complètement décorrélées.
		\pause
		\item Qu'une sélection des $k$ plus grandes composantes principales permet de \textbf{conserver} une grande partie de la \textbf{variance} tout en \textbf{réduisant} la \textbf{dimensionnalité} du problème.
		\pause
		\item Comment un algorithme d'apprentissage automatique utilise des images comme données d'entrée.
		\pause
		\item Comment l'analyse en composantes principales a réussi à simplifier les images d'un dataset de chiffres écrits à la main en divisant par 5 la taille des données tout en préservant 95 \% de la variance.
	\end{itemize}
\end{frame}